\documentclass[12pt]{report}
\usepackage{polski}
\usepackage[a4paper, margin=2.5cm]{geometry}
\usepackage{graphics}
\usepackage[export]{adjustbox}
\usepackage{float}
\usepackage{amsmath}
\usepackage{xcolor}
\usepackage{multirow}
\usepackage{array}

\begin{document}
	\section{Format pliku wejściowego}
	\subsection{Format binarny}
	\paragraph{}
	Pliki z danymi pomiarowymi generowane przez maszynę wytrzymałościową z rozszerzeniem \textit{.W01}, zapisane są w niestandardowym formacie binarnym. Aby zczytać dane z plików i stworzyć algorytm importujący do programu, przeprowadzona została analiza formatu. Plik zawiera nagłówek z metadanymi i informacjami o rozmiarze oraz zmienną ilość serii danych o zmiennej długości.
	
	\subsection{Struktura pliku}
	
	\subsubsection{Nagłówek}
	\paragraph{}
	W nagłówgu pliku, który znajduje się na początku, przed właściwymi danymi liczbowymi znajdują się metadane, informacje o rozmiarze, godzina utworzenia oraz nazwa programu wejściowego do maszyny badawczej. Cały nagłówek ma stały rozmiar 208 bajtów.
	
	\begin{table}[H]
		\centering
		\caption{Struktura nagłówka pliku}
		\begin{tabular}{|c|c|c|c|p{5.5cm}|} \hline
			offset & rozmiar & wartość/typ & nazwa & komentarz \\\hline
			\hline
			0 &  4 & 0x3f866666 & magic & stała wartość charakteryzująca format \\\hline
			4 &  4 & 1 & nieznane & najpewniej wersja formatu \\\hline
			8 &  4 & $rozmiar\_pliku + 1$ & rozmiar & \\\hline
			12 &  4 & 0 & nieznana &  \\\hline
			16 &  8 & "ScriptPl" & stały napis & ciąg znaków bez \verb|\0| \\\hline
			24 &  4 & $rozmiar\_pliku - 16$ & rozmiar2 &  \\\hline
			28 &  4 & liczba wierszy & liczba wierszy & ilość liczb w każdej kolumnie danych \\\hline
			32 &  4 & liczba kolumn & liczba kolumn & ilość kolumn z danymi \\\hline
			36 &  4 & 1 & nieznane & \\\hline
			40 &  4 & 1 & nieznane & \\\hline
			44 & 18 & "dd-MM-yyyyhh:mm:ss" & czas utworzenia & ciąg znaków bez \verb|\0| \\\hline
			62 & 42 & nieznana & nieznane & 42 bajty, głównie zera \\\hline
			104 & 104 & nazwa pliku & plik badania & nazwa pliku wejściowego do maszyny, dopełniona zerami do 104 bajtów \\\hline
		\end{tabular}
	\end{table}
	
	\pagebreak
	\subsubsection{Kolumny (serie danych)}
	\paragraph{}
	Kolumny mają zmienny rozmiar, zależnie od pliku, jednak zawsze taki sam w obrębie jednego pliku.
	
	\begin{table}[H]
		\centering
		\caption{Struktura kolumny}
		\begin{tabular}{|c|c|c|c|p{6.5cm}|} \hline
			offset & rozmiar & wartość/typ & nazwa & komentarz \\\hline
			\hline
			0 & 4 & int & nieznane & 1 lub 16385  \\\hline
			4 & 59 & ciąg znaków & nazwa & nazwa i jednostka kolumny, oddzielone spacjami \\\hline
			63 & $4 * liczba\_wierszy$ & float & dane & tablica 32-bitowych liczb zmiennoprzecinkowych \\\hline
		\end{tabular}
	\end{table}
	
	\subsubsection{Struktura całego pliku}
	\paragraph{}
	Plik składa się z nagłówka o stałym rozmiarze i zmiennej liczby kolumn zmiennej długości, które zawierają swój nagłówek oraz danych liczbowych.
	
	\begin{table}[H]
		\centering
		\caption{Struktura pliku}
		\begin{tabular}{|p{5cm}|p{4cm}|} \hline
			rozmiar  & element \\\hline
			\hline
			208 & nagłówek  \\\hline
			$63+4*liczba\_wierszy$  & kolumna 1 \\\hline
			$63+4*liczba\_wierszy$  & kolumna 2 \\\hline
			...  & ... \\\hline
			$63+4*liczba\_wierszy$  & kolumna $n$ \\\hline
		\end{tabular}
	\end{table}
	
\end{document}